\documentclass[12pt,a4paper]{ctexart}

\usepackage[margin=2.5cm]{geometry}
\usepackage{setspace}
\usepackage{hyperref}
\usepackage{enumitem}
\usepackage{titlesec}

\setstretch{1.25}
\hypersetup{hidelinks}
\setlist[itemize]{leftmargin=2em}

\titleformat{\section}{\bfseries\zihao{4}}{\thesection}{0.75em}{}
\titleformat{\subsection}{\bfseries\zihao{-4}}{\thesubsection}{0.75em}{}

\title{\texorpdfstring{Sb$_2$Te$_3$}{Sb2Te3}基异质结的性能调控研究进展\\\large 综述论文选题名称与大纲}
\author{}
\date{2026年5月}

\begin{document}

\maketitle

\section*{选题说明}
本文选定的综述题目为“\textbf{Sb$_2$Te$_3$基异质结的性能调控研究进展}”，聚焦于Sb$_2$Te$_3$与半导体、二维材料及相关功能层构成异质结后的能带、界面、载流子输运、热电、光电及拓扑关联性能的调控机制，不涉及“可控制备”或“应用”作为独立综述主题。Sb$_2$Te$_3$不仅是典型层状范德华材料，还兼具拓扑绝缘体、窄带隙半导体和热电材料等多重属性，因此非常适合作为异质集成中的功能核心层，用于实现电荷输运、光吸收、热电转换及自旋相关效应之间的协同调控[1-3]。

近年来，Sb$_2$Te$_3$基异质结在热电增强、宽谱光探测、范德华耦合调制以及缺陷诱导输运优化等方向均取得了新的进展[4-9]。据本文当前纳入的9篇代表性文献统计，发表于近三年（2023--2026年）的文献共5篇，占比约55.6\%，满足“近三年文献不少于三分之一”的要求[5-9]。

\section*{论文大纲}

\section{绪论}

\subsection{\texorpdfstring{Sb$_2$Te$_3$}{Sb2Te3}材料的基本性质与晶体结构}
Sb$_2$Te$_3$属于V--VI族层状窄禁带化合物，具有由五原子层（quintuple layer）堆垛形成的菱方结构，其层间以范德华作用结合，便于构筑低维和层状异质结[1,3]。同时，Sb$_2$Te$_3$还是Bi$_2$Se$_3$家族的重要拓扑绝缘体，其表面Dirac态与体态输运的耦合为性能调控研究提供了独特物理背景[1]。

\subsection{异质结的基本概念与性能调控意义}
异质结通常按能带对齐关系分为I型、II型和III型，不同对齐方式会显著影响界面电荷转移、势垒高度和载流子分离效率。对于Sb$_2$Te$_3$基体系，通过异质集成引入能带工程、界面散射和内建电场，可同时优化电输运、热输运和光电响应，是提升综合性能的核心路径[2,4-8]。

\subsection{本文综述范围与结构安排}
本文围绕“性能调控”这一主线展开，重点讨论Sb$_2$Te$_3$基异质结在能带结构、界面特征、载流子输运以及热电/光电相关性能方面的研究进展，并兼顾拓扑自旋相关行为。对于GeTe、磁性材料、超导材料等构筑的界面相变或近邻耦合体系，本文将其视为值得关注的拓展方向，在展望部分进行概括，而不单独扩展为主体章节。

\section{\texorpdfstring{Sb$_2$Te$_3$}{Sb2Te3}基异质结的能带结构调控}

\subsection{能带对齐类型及其对载流子输运的影响}
Sb$_2$Te$_3$与MoS$_2$、石墨烯、Te、MoTe$_2$等材料形成异质结后，能带匹配关系会决定界面势垒、载流子注入方向及复合概率，从而影响整流、光生分离和热电能量过滤行为[4-8]。

\subsection{组分调控对能带结构的调制}
通过掺杂、合金化，可以改变Sb$_2$Te$_3$附近费米能级位置、带边曲率及局域态分布，从而实现对载流子类型、浓度与有效质量的协同调控。AgSbTe$_2$--Sb$_2$Te$_3$异相合金的研究表明，组分无序和缺陷诱导的态密度重构能够显著改变热电输运特性[8，14-15]。

\subsection{应力工程与厚度效应对能隙的影响}
Sb$_2$Te$_3$层状结构对外部应变和厚度变化较为敏感[13]，晶格失配、柔性弯曲或层数变化都可能引起能隙、带边位置及表面态耦合的改变。针对Sb$_2$Te$_3$/MoS$_2$柔性异质结的研究表明，外部应变可直接调节其带结构和宽谱光探测行为[5]。

\section{\texorpdfstring{Sb$_2$Te$_3$}{Sb2Te3}基异质结的界面工程与性能调控}

\subsection{界面电荷转移与内建电场的形成}
界面功函数差和电负性差异会驱动电荷重新分布，进而形成内建电场，这一过程直接决定光生载流子的分离效率与跨界面输运能力。以1T$'$-MoTe$_2$/Sb$_2$Te$_3$超晶格样薄膜[12]为例，较弱层间相互作用与小功函数差共同影响了带弯曲和空穴输运[6]。

\subsection{界面缺陷态对电学与热学性能的影响}
空位、反位和无序界面既可能引入额外散射中心、降低迁移率，也可能通过能量过滤或抑制声子传输改善热电性能，因此其作用具有“双刃剑”特征。相关研究表明，适度缺陷和异相界面可有效降低热导率并重塑费米能级附近态密度，是Sb$_2$Te$_3$基异质结构性能优化的重要抓手[7-9]。

\subsection{范德华异质结中的层间耦合调控}
由于Sb$_2$Te$_3$可通过范德华作用与二维材料无悬挂键集成，层间耦合强弱成为调控电子结构和界面态分布的关键因素。Sb$_2$Te$_3$/graphene体系中的莫尔调制现象说明，即便在弱耦合界面中，周期势场仍可显著改变局域电子结构，为精细调控提供新思路[4]。

\section{\texorpdfstring{Sb$_2$Te$_3$}{Sb2Te3}基异质结的载流子浓度与输运性能调控}

\subsection{元素掺杂与异相复合对载流子浓度的调控}
元素掺杂、异相复合与化学计量比调控可通过改变材料的本征缺陷平衡和重构能带结构，实现载流子浓度、迁移率及有效质量的精准调节。不同掺杂元素和复合体系可针对性地解决单一材料载流子浓度过高或过低的问题，使电输运性能达到最优状态。

\subsection{外场调控对输运行为的影响}
电场、应力场等外场可作为非侵入式的调控手段，实时改变异质结的界面势垒高度和载流子空间分布。通过外场作用可灵活调节器件的输运特性，为柔性电子器件和可调控光电器件的设计提供了新的思路。

\subsection{界面散射与能量过滤效应}
异质界面、晶界与多尺度缺陷可同时对电子和声子产生不同程度的散射作用，适当的界面势垒能够有效筛除低能载流子并抑制声子传输。通过设计不同维度和结构的界面，可实现电输运与热输运的解耦调控，为功能材料的性能优化提供了核心途径。

\section{\texorpdfstring{Sb$_2$Te$_3$}{Sb2Te3}基异质结的热电与光电性能调控}

\subsection{热电性能提升的关键路径}
Sb$_2$Te$_3$基热电材料的性能优化始终围绕提高功率因子和降低热导率两条核心主线展开，通过载流子调控、界面工程与微观结构设计的协同作用实现热电转换效率的提升。从体材料到低维薄膜、多层膜及量子阱结构的发展，进一步拓展了热电性能的优化空间。

\subsection{光电响应性能的调控机制}
Sb$_2$Te$_3$基材料凭借窄带隙和高载流子迁移率的本征优势，通过构建异质结可优化光吸收、载流子分离及收集等光电转换的关键环节。异质结策略可显著提升器件的光电响应性能，并拓展其在宽谱探测、太阳能转换及超快光子学等领域的应用。

\subsection{多性能耦合调控的研究趋势}
未来Sb$_2$Te$_3$基材料的研究将从单一性能的提升转向电、热、光、自旋及相变特性的多维度耦合调控。通过挖掘材料的多重本征特性，可实现多功能集成器件的设计与制备，推动新一代低功耗、高性能电子系统的发展。

\section{\texorpdfstring{Sb$_2$Te$_3$}{Sb2Te3}基异质结的拓扑自旋相关性能调控}

\subsection{拓扑表面态与自旋动量锁定特征}
Sb$_2$Te$_3$作为典型三维拓扑绝缘体，其表面态具有自旋动量锁定特征，这使异质结中的输运行为不再完全服从传统半导体界面模型。该部分重点讨论拓扑表面态在界面耦合、自旋极化输运及表面散射抑制中的基础作用，并说明其为何构成Sb$_2$Te$_3$基异质结区别于普通热电或光电异质结的重要物理特征[1,4]。

\subsection{异质界面对拓扑自旋输运的调节机制}
当Sb$_2$Te$_3$与石墨烯、二维过渡金属硫属化物或金属功能层形成异质界面时，界面电势、轨道杂化与局域对称性破缺都会改变表面态分布及自旋输运通道。已有研究表明，范德华耦合、莫尔调制和界面散射重构可影响表面态保持程度与自旋相关输运响应，因此拓扑自旋效应应作为独立章节加以讨论[4-6]。

\subsection{拓扑自旋性能调控的研究意义与发展趋势}
拓扑自旋相关研究不仅有助于深化对Sb$_2$Te$_3$基异质结界面物理的理解，也为低功耗自旋电子学和多功能耦合器件提供了理论基础。未来可围绕高质量界面构筑、自旋分辨表征、磁邻近效应引入及拓扑态与热电/光电性能协同优化展开更系统研究；若与磁性材料或超导材料耦合，还可能进一步拓展到交换场调控、超导近邻效应与量子输运等更前沿的性能维度。

\section{总结与展望}

\subsection{现有研究的阶段性认识}
现有研究表明，Sb$_2$Te$_3$基异质结的性能调控已从单纯的材料组合逐步发展为围绕能带匹配、界面耦合、缺陷调控和多物理场响应的系统设计。不同调控路径之间并非彼此孤立，而是共同作用于载流子输运、声子散射、界面电势分布、界面相变行为以及拓扑自旋相关输运等核心物理过程[4-11]。

\subsection{面临的挑战与未来发展方向}
后续研究仍需解决界面结构表征不足、调控机制定量化不充分以及多性能指标难以兼顾等问题。未来可重点关注原子级界面设计、原位表征结合理论计算、面向低功耗器件的性能评价体系，以及近室温条件下热电—光电—拓扑耦合机制的统一认识。除此之外，Sb$_2$Te$_3$与GeTe、磁性材料和超导材料构成的异质结构还提示出界面相变耦合、磁近邻调控和超导近邻效应等潜在写作角度，但考虑到这部分文献目前相对分散，更适合在综述中作为前沿展望而非主体章节展开。

\section*{参考文献}
\begin{thebibliography}{99}
\bibitem{ref1} Zhang H, Liu C X, Qi X L, et al. Topological insulators in Bi$_2$Se$_3$, Bi$_2$Te$_3$ and Sb$_2$Te$_3$ with a single Dirac cone on the surface[J]. Nature Physics, 2009, 5(6): 438-442.
\bibitem{ref2} Venkatasubramanian R, Siivola E, Colpitts T, et al. Thin-film thermoelectric devices with high room-temperature figures of merit[J]. Nature, 2001, 413(6856): 597-602.
\bibitem{ref3} Peranio N, Winkler M, Bessas D, et al. Room-temperature MBE deposition, thermoelectric properties, and advanced structural characterization of binary Bi$_2$Te$_3$ and Sb$_2$Te$_3$ thin films[J]. Journal of Alloys and Compounds, 2012, 521: 163-173.
\bibitem{ref4} Yin Y, Wang G, Liu C, et al. Moir\'e-pattern-modulated electronic structures in Sb$_2$Te$_3$/graphene heterostructure[J]. Nano Research, 2022, 15(2): 1115-1119.
\bibitem{ref5} Wang H, Dong C, Gui Y, et al. Self-powered Sb$_2$Te$_3$/MoS$_2$ heterojunction broadband photodetector on flexible substrate from visible to near infrared[J]. Nanomaterials, 2023, 13(13): 1973.
\bibitem{ref6} Li Z, Zhang C, Luo J, et al. Weak interlayer interactions and nearly temperature independent electrical transport in p-type 1T$'$-MoTe$_2$/Sb$_2$Te$_3$ superlattice-like films[J]. Journal of Solid State Chemistry, 2024, 336: 124785.
\bibitem{ref7} Shahid I, Zhang X, Ali A, et al. Theoretical insights into Sb$_2$Te$_3$/Te van der Waals heterostructures for achieving very high figure of merit and conversion efficiency[J]. International Journal of Heat and Mass Transfer, 2025, 238: 126479.
\bibitem{ref8} Kaur R, Tanwar A, Padmanathan N, et al. Anomalous thermoelectric nature in disordered AgSbTe$_2$-Sb$_2$Te$_3$ hetero-phase alloys for room temperature applications[J]. Journal of Alloys and Compounds, 2025, 1010: 177313.
\bibitem{ref9} Mori R, Kurokawa T, Yamauchi K, et al. Improved thermoelectric performances of nanocrystalline Sb$_2$Te$_3$/Cr bilayers by reducing thermal conductivity in the grain boundaries and heterostructure interface[J]. Vacuum, 2019, 161: 92-97.
\bibitem{ref10} Ferhat M, Zaoui A, Certier M, et al. Structural, electronic and optical properties of GeTe/Sb$_2$Te$_3$ superlattices: a first-principles study[J]. Superlattices and Microstructures, 2018, 121: 139-146.
\bibitem{ref11} Mogi M, Kawamura M, Yoshimi R, et al. Tailoring tricolor structure of magnetic topological insulator for robust axion insulator[J]. Science Advances, 2017, 3(4): eaao1669.
\bibitem{ref12} 唐容喆,胡松柏,高静静,等. Sb$_2$Te$_3$薄膜的性质及光伏应用[J]. 太阳能, 2012(23): 3.
\bibitem{ref13} 赵堃. 压力下Bi$_2$Te$_3$和Sb$_2$Te$_3$中电子拓扑转变的第一性原理研究[D]. 哈尔滨: 哈尔滨工业大学, 2016.
\bibitem{ref14} 王汉文. 范德华层状材料Sb$_2$Te$_3$和Bi$_2$Te$_3$中3d过渡金属各向异性扩散机制[D]. 武汉: 武汉理工大学, 2024.
\bibitem{ref15} 谢亮军. Sb$_2$Te$_3$基半导体的组织结构与热电性能[D]. 哈尔滨: 哈尔滨工业大学, 2019.
\bibitem{ref16} 王驰,王钟. 基于电子输运调控的Ti掺杂Sb$_2$Te$_3$半导体工艺分析[J]. 集成电路应用, 2025, 42(05): 397-399.
\bibitem{ref17} 张梦雨. 高热稳定性C掺杂Sb$_2$Te$_3$相变材料研究[D]. 上海: 东华大学, 2024.
\bibitem{ref18} 易文,赵永杰,王伯宇,等. Sb$_2$Te$_3$基热电薄膜的研究进展[J]. 硅酸盐学报, 2021, 49(06): 1111-1124.
\bibitem{ref19} 张度. 基于窄带隙二维Sb$_2$Te$_3$的宽光谱偏振光电探测器[D]. 北京: 北京信息科技大学, 2025.
\bibitem{ref20} 缪家乐. GeTe/Sb$_2$Te$_3$纳米多层膜及量子阱热电性能研究[D]. 镇江: 江苏科技大学, 2023.
\bibitem{ref21} 张优. 成分与微观结构调制对P型（Bi,Sb）$_2$Te$_3$合金热电性能的优化[D]. 秦皇岛: 燕山大学, 2022.
\end{thebibliography}

\end{document}
